% Options for packages loaded elsewhere
\PassOptionsToPackage{unicode}{hyperref}
\PassOptionsToPackage{hyphens}{url}
\documentclass[
]{article}
\usepackage{xcolor}
\usepackage[margin=1in]{geometry}
\usepackage{amsmath,amssymb}
\setcounter{secnumdepth}{5}
\usepackage{iftex}
\ifPDFTeX
  \usepackage[T1]{fontenc}
  \usepackage[utf8]{inputenc}
  \usepackage{textcomp} % provide euro and other symbols
\else % if luatex or xetex
  \usepackage{unicode-math} % this also loads fontspec
  \defaultfontfeatures{Scale=MatchLowercase}
  \defaultfontfeatures[\rmfamily]{Ligatures=TeX,Scale=1}
\fi
\usepackage{lmodern}
\ifPDFTeX\else
  % xetex/luatex font selection
\fi
% Use upquote if available, for straight quotes in verbatim environments
\IfFileExists{upquote.sty}{\usepackage{upquote}}{}
\IfFileExists{microtype.sty}{% use microtype if available
  \usepackage[]{microtype}
  \UseMicrotypeSet[protrusion]{basicmath} % disable protrusion for tt fonts
}{}
\makeatletter
\@ifundefined{KOMAClassName}{% if non-KOMA class
  \IfFileExists{parskip.sty}{%
    \usepackage{parskip}
  }{% else
    \setlength{\parindent}{0pt}
    \setlength{\parskip}{6pt plus 2pt minus 1pt}}
}{% if KOMA class
  \KOMAoptions{parskip=half}}
\makeatother
\usepackage{longtable,booktabs,array}
\newcounter{none} % for unnumbered tables
\usepackage{calc} % for calculating minipage widths
% Correct order of tables after \paragraph or \subparagraph
\usepackage{etoolbox}
\makeatletter
\patchcmd\longtable{\par}{\if@noskipsec\mbox{}\fi\par}{}{}
\makeatother
% Allow footnotes in longtable head/foot
\IfFileExists{footnotehyper.sty}{\usepackage{footnotehyper}}{\usepackage{footnote}}
\makesavenoteenv{longtable}
\usepackage{graphicx}
\makeatletter
\newsavebox\pandoc@box
\newcommand*\pandocbounded[1]{% scales image to fit in text height/width
  \sbox\pandoc@box{#1}%
  \Gscale@div\@tempa{\textheight}{\dimexpr\ht\pandoc@box+\dp\pandoc@box\relax}%
  \Gscale@div\@tempb{\linewidth}{\wd\pandoc@box}%
  \ifdim\@tempb\p@<\@tempa\p@\let\@tempa\@tempb\fi% select the smaller of both
  \ifdim\@tempa\p@<\p@\scalebox{\@tempa}{\usebox\pandoc@box}%
  \else\usebox{\pandoc@box}%
  \fi%
}
% Set default figure placement to htbp
\def\fps@figure{htbp}
\makeatother
\setlength{\emergencystretch}{3em} % prevent overfull lines
\providecommand{\tightlist}{%
  \setlength{\itemsep}{0pt}\setlength{\parskip}{0pt}}
\usepackage{bookmark}
\IfFileExists{xurl.sty}{\usepackage{xurl}}{} % add URL line breaks if available
\urlstyle{same}
\hypersetup{
  pdftitle={Latent-class and mixture-model formulations for biomarker-treatment interaction in N-of-1 trials},
  pdfauthor={pmsimstats team},
  hidelinks,
  pdfcreator={LaTeX via pandoc}}

\title{Latent-class and mixture-model formulations for biomarker-treatment interaction in N-of-1 trials}
\author{pmsimstats team}
\date{2026-05-06}

\begin{document}
\maketitle

\section{Abstract}\label{abstract}

\emph{Draft placeholder.} Architectures A and B in the companion paper
encode the biomarker-treatment interaction as either a deterministic
moderation of the mean (A) or a covariance-structure correlation that
erodes under carryover (B). A reviewer of that paper raised a
substantive question: if Architecture B is functionally an imperfect
responder indicator, why is the heterogeneity not modelled
explicitly, at the participant level, via a latent-class or
mixture-model construction? This paper takes that suggestion as its
starting point. We develop the latent-class and finite-mixture
formulations of the biomarker-treatment interaction problem,
characterise their identifiability and estimability under the
N-of-1 family of trial designs, and contrast their statistical
properties against Architectures A and B as a function of trial
length, biomarker measurement error, carryover half-life, and
between-participant heterogeneity. The full source for the
exposition that motivated this paper is preserved in
\texttt{notes/source-revision-latent-class.md}.

\section{Introduction}\label{introduction}

\emph{Section stub.} Frame the question in terms of the dual-tradition
gap identified in the source revision response: regression-with-
random-effects (biostatistical default) versus latent-variable
construction (psychometric and econometric default). State the
contribution: first principled application of latent-variable
machinery to the biomarker-treatment interaction problem under the
N-of-1 design family.

\section{Statistical formulation}\label{statistical-formulation}

\emph{Section stub.} Develop:

\begin{enumerate}
\def\labelenumi{\arabic{enumi}.}
\tightlist
\item
  A two-class latent-class model with class membership predicted by
  the biomarker.
\item
  A finite-mixture-of-regressions model with biomarker-dependent
  mixing weights.
\item
  Connections to the random-coefficient growth-mixture-model
  tradition.
\item
  Identifiability conditions specific to the N-of-1 setting (within-
  participant treatment switches as a key source of identification).
\end{enumerate}

\section{Estimation}\label{estimation}

\emph{Section stub.} EM algorithms, MCMC samplers, and the comparative
behaviour of point versus posterior estimates. Software:
\texttt{mclust}, \texttt{flexmix}, \texttt{OpenMx}, \texttt{Stan} candidate implementations.

\section{Simulation study}\label{simulation-study}

\emph{Section stub.} Re-use the trial-design grid and biomarker
parameter space from the carryover-sensitivity factorial. Add a
dimension over latent-class separability. Outcomes: power for the
class-by-treatment interaction, classification accuracy, recovery
of true class means.

\section{Application}\label{application}

\emph{Section stub.} Re-analyse the prazosin-PTSD data with the latent-
class formulation; compare to the Architecture A and B fits in the
companion papers.

\section{Discussion}\label{discussion}

\emph{Section stub.} When does the latent-class construction add value
over Architecture A or B? When does the data not support its
identifiability?

\section{Acknowledgements}\label{acknowledgements}

\section{References}\label{references}

\end{document}
